% chapters/04/evaluasi.tex
\bagian{Asesmen dan Instrumen Penilaian}
\label{sec:asesmen-game-loop}

\begin{asesmenbox}
\textbf{Kedudukan asesmen bab ini.} Bab ini menuntaskan \SUBCPMK{1.1} secara
utuh, sehingga ketiga indikatornya diperiksa. Asesmennya terdiri atas dua
bagian: penilaian artefak program melalui Tugas Praktikum 02, dan soal konsep
yang menuntut penjelasan.

\vspace{0.4em}
Selain dinilai melalui kedua instrumen di bawah ini, materi Bab~IV termasuk
bahan \textbf{Kuis 1} pada minggu kelima (Materi 1--4) dan bahan
\textbf{Ujian Tengah Semester} pada minggu kesepuluh yang mengukur
\CPMK{1} dan \CPMK{2}.
\end{asesmenbox}

\subsection*{Instrumen 1: Rubrik Praktikum 02}

\begin{table}[htbp]
\centering
\caption{Rubrik penilaian Praktikum 02 beserta indikator yang diukur}
\label{tab:rubrik-praktikum-02}
{\setlength{\tabcolsep}{4pt}\footnotesize%
\begin{tabularx}{\textwidth}{@{}L{3.2cm} C{1.3cm} C{2.1cm} Y@{}}
\toprule
\headrow \textbf{Kriteria} & \textbf{Bobot} & \textbf{Indikator} & \textbf{Deskripsi capaian} \\
\midrule
Siklus berjalan dan dapat diakhiri & 35\% & \ind{a} &
Ketiga tahap dikerjakan berulang kali, jendela menanggapi masukan tanpa
menggantung, dan program berakhir dengan rapi melalui tombol maupun penutupan
jendela. \\
\addlinespace
Pengendalian laju dan pengukurannya & 35\% & \ind{b} &
Laju bingkai dibatasi dengan \code{clock.tick()}, laju terukur ditampilkan
pada judul jendela, dan hasil pengukuran pengaruh batas laju terhadap
pemakaian prosesor dicatat pada lembar kerja. \\
\addlinespace
Kerapian kode dan dokumentasi & 30\% & \ind{a}, \ind{c} &
Ketiga tahap terlihat jelas pada kode, keadaan dibuat di luar perulangan,
terdapat komentar pada bagian penting, dan catatan praktikum memuat hasil
pengukuran beserta kesimpulannya. \\
\bottomrule
\end{tabularx}}
\end{table}

\subsection*{Instrumen 2: Soal Konsep}

Jawablah dengan singkat, tanpa membuka bab terlebih dahulu. Kunci jawaban
disajikan pada Tabel~\ref{tab:kunci-konsep-ch04}.

\begin{enumerate}
  \item Jelaskan urutan ketiga tahap siklus permainan beserta pekerjaan pada
        masing-masing tahap.
  \item Apa dua pekerjaan yang dilakukan \code{clock.tick(60)}?
  \item Mengapa \code{pygame.event.wait()} tidak dapat dipakai di dalam siklus
        permainan?
  \item Jelaskan perbedaan langkah tetap dan langkah berubah beserta satu
        kelebihan masing-masing.
  \item Mengapa menaikkan batas laju tidak selalu menaikkan laju yang
        tercapai?
\end{enumerate}

\begin{table}[htbp]
\centering
\caption{Kunci soal konsep Bab IV beserta indikator yang diukur}
\label{tab:kunci-konsep-ch04}
{\setlength{\tabcolsep}{4pt}\footnotesize%
\begin{tabularx}{\textwidth}{@{}C{0.8cm} C{2.1cm} Y@{}}
\toprule
\headrow \textbf{No} & \textbf{Indikator} & \textbf{Pokok jawaban} \\
\midrule
1 & \ind{a} &
Masukan membaca peristiwa pemain; pembaruan mengubah keadaan permainan
berdasarkan masukan dan waktu; penayangan menggambar ulang seluruh jendela
menurut keadaan terbaru. \\
\addlinespace
2 & \ind{b} &
Mengukur selang waktu yang benar-benar berlangsung sejak pemanggilan
sebelumnya, dan menahan laju dengan menunggu sampai jatah waktu satu bingkai
terpenuhi. \\
\addlinespace
3 & \ind{a} &
Karena \code{event.wait()} menghentikan program sampai ada peristiwa. Ketika
pemain diam, tidak ada peristiwa yang datang, sehingga tahap pembaruan dan
penayangan tidak pernah dikerjakan dan permainan membeku. \\
\addlinespace
4 & \ind{c} &
Langkah tetap memakai panjang langkah yang sama sehingga hasilnya dapat diulang
dan perhitungan fisika lebih stabil; langkah berubah mengikuti lama bingkai
sesungguhnya sehingga lebih sederhana diterapkan. \\
\addlinespace
5 & \ind{b} &
Karena laju yang tercapai ditentukan oleh lama pekerjaan satu bingkai. Apabila
pekerjaan itu lebih lama daripada jatah waktu pada batas yang diminta, siklus
berjalan secepat kemampuan program dan batasnya tidak pernah tercapai. \\
\bottomrule
\end{tabularx}}
\end{table}
